\chapter{Transformer 核心机制：Attention、FFN、Norm 与损失}

\begin{introduction}
\item {\heiti 本章核心问题}：模型为什么能够基于上下文继续生成，而企业知识助手中常见的“遗忘、跑偏、格式漂移、训练不稳”等现象，又分别对应哪些 Transformer 核心机制？
\item {\heiti 主案例推进到哪一步}：把企业知识助手中的典型错误症状，映射回 Attention、FFN、Norm、位置设计和训练目标这几类底层部件。
\item {\heiti 读完后能做什么}：看到长上下文答偏、微调不稳、格式失控或训练目标错配时，能先回到正确的底层机制解释它，而不是只会说“模型太笨”。
\end{introduction}

\section{学习目标}
\begin{itemize}
  \item 理解 Transformer 中最关键的几类部件：Attention、FFN、Norm、位置设计与训练目标。
  \item 从单层数据流的角度看清模型在每一步到底做了什么，而不是把各个名词拆开死记。
  \item 能把企业知识助手中的真实工程症状，映射回相应的底层机制。
  \item 建立一个最低限度但足够可靠的原理框架，为后面的推理、微调、RAG 和性能优化打底。
\end{itemize}

\section{问题背景}
很多工程问题表面上看像是“调参问题”，本质上却来自对底层机制的误判。比如，为什么长上下文时模型更容易答偏？为什么同样是微调，有的模型很稳、有的模型频繁 loss spike？为什么同样是生成任务，有的模型格式遵循更好、有的模型却总在关键字段上漂移？这些问题如果只靠经验记忆，很难迁移到新的模型和新的框架中。

因此，本章不追求完整推导，而是只保留工程读者必须掌握的一条数据流叙事：\textbf{模型先把 token 变成向量，借助位置设计知道“谁在前谁在后”，再通过 Attention 决定看哪里，通过 FFN 决定怎样改写表示，在 Norm 和残差路径的帮助下稳定传播，最后在交叉熵或 KL 这样的目标下被训练。}

这条叙事的价值在于，后面无论你是在看推理链路、做微调、调长上下文、搭 RAG，还是分析训练不稳定，最终都会不断回到这条主线。如果第一遍就把这条主线看清，后面的很多章节会轻松很多。

\section{核心概念}
\subsection{先把一个 Transformer Block 看成一条数据流}
学习 Transformer 最容易陷入的误区，是把 Attention、FFN、Norm、RoPE、损失函数分别当作独立知识点。更稳的方式是先把它看成一条连续数据流。对每个 token 来说，一层 Transformer 至少会经历下面几步：
\begin{enumerate}
  \item 输入 token 被映射成向量表示。
  \item 位置设计告诉模型这个 token 在序列里的相对位置。
  \item Attention 决定它应该从上下文中的哪些位置吸收信息。
  \item FFN 决定吸收完信息后，这个位置的表示应如何被非线性改写。
  \item Norm 和残差路径保证整个深层网络训练和传播时不至于失稳。
  \item 最终顶层输出一个词表分布，由训练目标决定它该往哪个方向继续优化。
\end{enumerate}

只要记住这六步，后面遇到绝大多数模型结构问题时，就不会完全失去方向。因为你至少能问自己三件事：现在的问题是“看错了地方”，还是“改写错了表示”，还是“训练目标把模型推向了不对的方向”？

\begin{figure}[H]
\centering
\begin{tikzpicture}[
  font=\small,
  box/.style={llm box, minimum width=2.9cm, minimum height=1cm},
  note/.style={llm note box, minimum width=2.8cm, minimum height=0.9cm},
  arr/.style={llm arr}
]
\node[box] (embed) {Token 表示\\嵌入与位置};
\node[box, right=0.8cm of embed] (attn) {Attention\\决定看哪里};
\node[box, right=0.8cm of attn] (ffn) {FFN\\决定怎样改写};
\node[box, right=0.8cm of ffn] (head) {输出分布\\预测下一个 token};
\node[note, below=0.9cm of attn] (norm) {Norm + 残差\\稳住深层传播};
\node[note, below=0.9cm of head] (loss) {交叉熵 / KL\\规定优化方向};
\draw[arr] (embed) -- (attn);
\draw[arr] (attn) -- (ffn);
\draw[arr] (ffn) -- (head);
\draw[arr] (norm) -- (attn);
\draw[arr] (norm) -- (ffn);
\draw[arr] (loss) -- (head);
\end{tikzpicture}
\caption{Transformer 单层的工程化数据流：先看哪里，再改写表示，最后在目标函数约束下学习}
\label{fig:transformer-block}
\end{figure}

\subsection{Token 表示与位置设计：模型先要知道“当前在看什么”}
对模型来说，原始文本并不是直接可计算的对象。文本首先会被切成 token，再映射成向量。这个向量不是“词义的唯一真相”，而是一个可训练、可更新的表示起点。模型后续所有阅读和推理，都是在这些向量空间里进行的。

但只有 token 向量还不够。因为序列模型不仅要知道“有哪些词”，还要知道“谁在前、谁在后、谁和谁距离近”。否则，“审批先经理后总监”和“审批先总监后经理”会被看成同一堆 token。位置设计的职责，就是让模型在向量空间里感知顺序和相对关系。

在现代 Decoder-Only 大模型里，最常见的位置设计是 RoPE 一类相对位置编码。为了帮助理解，先区分两种思路：

\begin{itemize}
  \item \textbf{绝对位置编码}（原始 Transformer 使用）：给每个位置分配一个固定向量，直接加到 token 嵌入上。缺点是训练时见过的最大长度就是它的上限，超出后表现急剧下降。
  \item \textbf{相对位置编码}（RoPE 属于此类）：不给每个位置一个固定标签，而是让模型在计算 Attention 时感知"两个 token 之间隔了多远"。好处是更容易泛化到训练时未见过的长度。
\end{itemize}

RoPE 的具体做法是对 Query 和 Key 向量施加一个与位置相关的旋转变换，使得两个位置的点积只取决于它们的相对距离。工程上你不需要先把复数旋转公式完全推清，只要先知道两件事：
\begin{itemize}
  \item 它的作用不是增加新知识，而是让 Attention 知道不同位置之间的相对关系。
  \item 长上下文能力不只是“把窗口拉长”，还要看位置设计在更长距离上是否还能稳定工作。
\end{itemize}

这也是为什么后面讲长上下文和 RAG 时，我们会反复回到“输入顺序、位置关系和片段排列”。模型并不是只读内容，它也在读位置。

\subsection{Tokenizer 决定模型到底把什么当成“一个单位”}
很多入门读者会把 tokenization 看成纯实现细节，仿佛只要模型够强，怎么切都无所谓。工程上这恰恰是一个很危险的误解。因为模型从来不是直接读“词”或“句”，它读的是 tokenizer 切出来的最小单位。也就是说，领域术语、错误码、表头字段、金额单位、代码标识符和中英文混排，首先都会在 tokenizer 这一关被重新定义。

对现代大模型来说，常见做法是 BPE、SentencePiece 这类子词切分。它们的好处是能用有限词表覆盖大量文本，但代价是：同一个业务实体在不同上下文里可能被切成不同粒度。例如“VPN-4037”可能被切成若干子块，“P6 级别”里的等级标识也可能和普通文字混在一起。只要这些切分在训练语料里不稳定，模型后面在 Attention 和输出词表阶段就会更难稳定抓住它们。

\begin{table}[H]
\centering
\small
\begin{tabularx}{\textwidth}{>{\raggedright\arraybackslash}p{2.8cm}>{\raggedright\arraybackslash}p{3cm}>{\raggedright\arraybackslash}X}
\toprule
\textbf{输入类型} & \textbf{tokenization 最容易出什么问题} & \textbf{为什么这会影响后续模型行为} \\
\midrule
领域术语、产品名 & 被切得过碎，稳定性差 & 模型更难把它们当成一个稳定语义对象去关注和输出 \\
错误码、版本号、表头字段 & 精确字符串被拆散 & 检索、重排序和结构化输出都更容易丢掉关键字面线索 \\
中英混排、代码片段 & 切分粒度忽粗忽细 & 上下文中相邻 token 的关系会更难学稳 \\
长数字与金额单位 & 数字与单位解绑 & 模型容易认识数值，却接错适用对象或计量语义 \\
\bottomrule
\end{tabularx}
\caption{Tokenizer 不只是把文本切开，它会决定模型后面学习的最小语义单位}
\label{tab:tokenizer-impact}
\end{table}

这也是为什么企业场景里经常会看到一种现象：模型好像“懂大意”，却总在产品型号、制度编号、错误码和字段名上不够稳。很多时候，这并不一定说明 Attention 没工作，而是更前面的 tokenizer 就已经把这些对象切得不够利于建模了。教材里保留这一层，是为了让读者知道：模型表示的第一步，往往就已经在决定后面哪些信息更容易被学会、哪些更容易一直发虚。

\subsection{输入嵌入与输出词表：模型从词表开始，也回到词表结束}
很多入门读者会把“输入 token 变成向量”理解成一个单向入口，好像模型只是在最开始用一次词表。其实对自回归语言模型来说，词表会在\textbf{起点和终点}同时出现。起点是输入嵌入矩阵，把 token id 变成隐藏表示；终点是输出头，把最终隐藏状态重新投影回整个词表的 logits，再决定下一个 token 的概率分布。

这条路径的工程含义很大。因为模型并不是在一个抽象空间里“想完了再随便写词”，而是每一步都要重新回到具体词表上做概率选择。也就是说，领域术语、产品型号、错误码、表头字段和格式标记，最后都要通过这个输出词表被真正选出来。若词表切分不理想、术语在训练数据里过少，或者某些结构化标记分布极不稳定，模型后面的生成就会更吃力。

不少现代模型还会把输入嵌入和输出投影做权重共享。这种做法的直觉是：\textbf{模型用什么空间去读 token，也希望大体在同一个空间里决定怎么把它们写出来。} 对工程读者来说，不必先陷进全部实现细节，但要先记住：输入嵌入、隐藏状态和输出词表并不是彼此断开的三段，而是一条前后闭环的数据流。

\subsection{位置编码扩展不是“把窗口拉长”这么简单}
旧稿里关于长文本处理讲过一个很重要的事实：\textbf{上下文窗口变长时，首先被考验的往往不是参数量，而是位置表示还能不能稳定工作。} 模型之所以能区分“标题在前、条件在后”和“条件在前、结论在后”，本质上依赖位置设计参与了表示构建。

一旦序列明显变长，工程上最常出现三类现象：
\begin{itemize}
  \item \textbf{局部关系还能读，远距离关系开始漂}：模型知道每句话在说什么，却很难把分散在远处的条件重新对齐。
  \item \textbf{检索片段一多就开始答偏}：不是片段没放进去，而是位置信号和噪声一起变复杂。
  \item \textbf{训练和推理成本同时上升}：位置扩展往往不是单独发生的，它会牵连 Attention、KV Cache 和吞吐。
\end{itemize}

因此，长文本问题从来不只是“窗口长度”问题，而是“位置信号、注意力成本和证据密度”共同作用的问题。这也是为什么真实系统里，长上下文和 RAG 往往不是二选一，而是互相分担压力。

\subsection{为什么条件和结论一旦隔太远，模型就更容易接错}
企业文档里经常会出现一种长距离依赖：前面先列适用条件，后面若干段以后才给出结论、例外或审批要求。对人来说，我们会一边读一边把条件暂存在脑中；对模型来说，这种跨距离绑定必须依赖位置表示和 Attention 在多层里持续把“谁约束谁”这件事保住。

一旦距离拉长、噪声增多，模型就更容易出现三类错误：把结论泛化到不适用的人群，把例外条件接到错误条款，或者只抓住主体结论却漏掉尾部限制。很多看上去像“模型理解力不够”的问题，本质上其实是长距离条件绑定在多层传播里逐渐变弱了。

这条机制对后面所有工程章节都很关键。它解释了为什么 RAG 不是把片段简单拼一起就够了，也解释了为什么表格、列表和制度条款一旦被粗暴切断，模型会尤其容易“每句话都认识，但整体关系接错”。当我们后面谈检索排序、上下文压缩和结构保留时，本质上都在想办法帮模型把这些长距离依赖重新压短、压清楚。

\subsection{Attention：模型如何在上下文中定位信息}
自注意力的核心计算通常写成
\[
\mathrm{Attention}(Q,K,V)=\mathrm{softmax}\left(\frac{QK^\top}{\sqrt{d_k}}\right)V.
\]
这里最重要的不是公式表面，而是它表达的行为：当前 token 会拿自己的 query 去和上下文里每个位置的 key 做匹配，再按匹配强度对 value 加权求和。也就是说，Attention 本质上是在做一种\textbf{上下文内检索}。

这个机制之所以关键，是因为语言中的有效信息通常并不只在当前位置附近。员工提问“住宿超标后还能不能报销”，模型也许需要同时看到前文中的“住宿标准”、后文中的“审批例外”和附件中的“需附审批记录”。Attention 让当前位置有能力跨越距离，从多个位置聚合相关信息。

因此，从工程角度看，Attention 解决的是“当前应从哪里取证”的问题。它不负责最终把结论写出来，但它决定了模型能不能先看对地方。只要这一步错了，后面再强的非线性变换也只是把错信息加工得更漂亮。

\subsection{Q、K、V 不是三个神秘字母：它们分别在问什么、标什么、带什么}
很多入门读者第一次看到 \(Q\)、\(K\)、\(V\) 时，会把它们当成公式里的三个抽象矩阵。更容易落地的理解是：\textbf{Query 在表达“我现在想找什么”，Key 在表达“我这里大概能提供什么”，Value 在携带“如果你真的来取，我会给你什么内容”。}

同一个 token 之所以要被投影成三种不同表示，是因为“我是谁”“我想找什么”“我能贡献什么信息”并不总是同一回事。比如在制度句子里，“审批”这个词作为 Query 时，可能是在寻找与自己相关的条件和角色；作为 Key 时，它是在告诉别的位置“我这里和流程、权限、条件有关”；作为 Value 时，它又携带了真正要被聚合走的内容表示。把这三件事拆开，模型就更容易在不同语境里灵活地建立关系。

这个区分还有一个很重要的工程含义：如果模型经常把看起来相近的字段、条款或角色混在一起，问题未必只是“没看到”，也可能是 Query 和 Key 的匹配空间没有把它们拉开得足够清楚。也就是说，Attention 的失败有时不是找不到信息，而是\textbf{把错误的信息判成了“像是我要找的那个”。} 后面我们讲重排序和检索对齐时，这个直觉会再次出现。

\subsection{Attention 不是事实库：它只会在当前可见上下文里搬运信息}
理解 Attention 时，一个特别值得尽早建立的边界是：\textbf{Attention 会在当前可见上下文里重新分配注意力，但它不会凭空制造当前上下文中不存在的证据。} 如果制度版本没被放进 prompt，或者表格字段在解析阶段已经丢了，那么 Attention 再强，也只能在剩下的材料里做最优搜索，而不是把缺失事实“找回来”。

这也是为什么很多人把模型答错误诊成“注意力机制不够先进”，其实真正缺的是知识接入、文档结构保留或版本过滤。Attention 能解释“为什么模型没看见已经在上下文里的那条关键信息”，却不能解释“为什么系统从来没把关键信息送进来”。前者是模型层问题，后者常常是系统层问题。

把这条边界记清楚，对后面读 RAG 和上下文管理非常有帮助。因为一旦上下文进入系统，你就要同时管两件事：一是让正确证据有机会进来，二是让模型在进来的证据里先看对地方。前者主要是检索与文档理解，后者才主要落在 Attention 这类内部机制上。

\subsection{因果 Mask：为什么生成模型只能看左边}
对 Decoder-Only 模型来说，还有一个经常被忽略但非常关键的部件：因果 Mask。它要求模型在预测当前 token 时，只能看左边已经出现的内容，不能偷看未来 token。没有这个限制，训练时模型会直接利用答案本身，生成任务就失去了意义。

这条约束决定了生成模型和双向理解模型的本质差别。双向模型更像“阅读”，为了理解某个词，可以同时看左右文；生成模型更像“写作”，只能基于已经写出的内容继续往后写。企业知识助手之所以围绕生成模型展开，也正是因为它的主要工作不是给现成句子打标签，而是把已有材料组织成新的回答。

\subsection{多头注意力：为什么不只看一种相关性}
如果只有一个注意力头，模型就只能用一种“相关性视角”看上下文。但真实语言里，相关性有很多种：有的头更容易盯住语法一致性，有的头更容易关注实体指代，有的头更偏向标题与正文的对应，有的头可能更擅长捕捉长距离约束。

多头注意力的意义，就是并行地提供多种检索视角，再把这些视角的结果拼起来。工程上你不用把每个头都神秘化，但要知道：多头不是为了“让图看起来更复杂”，而是为了让模型在同一层里能同时关心不同类型的依赖关系。

这对企业知识助手很重要。因为制度文档里的相关性往往并不单一。员工的一个问题，可能同时涉及主题匹配、条件限制、适用对象和例外条款。单一路径很容易漏，多个注意力头则更有机会把不同信号都带回来。

\subsection{Attention 为什么会成为工程瓶颈}
Attention 不只决定模型“看哪里”，还决定模型“为看这些地方付出多大代价”。随着上下文长度增加，Attention 计算和内存开销都会迅速上升，于是两个工程问题同时出现：
\begin{itemize}
  \item \textbf{训练侧}：显存占用增加，长上下文训练吞吐下降，更容易不稳定。
  \item \textbf{推理侧}：KV Cache 变大，解码时内存带宽压力上升，延迟更难压住。
\end{itemize}

这就是为什么后面讲长文档、RAG 和推理优化时，会反复回到 Attention。很多表面上像“部署问题”的症状，其实在这一层就已经埋下了根。长上下文并不是白来的能力，它总要用计算、显存和排序压力去换。

\begin{table}[H]
\centering
\small
\begin{tabularx}{\textwidth}{>{\raggedright\arraybackslash}p{2.9cm}>{\raggedright\arraybackslash}p{3cm}>{\raggedright\arraybackslash}X}
\toprule
\textbf{现象} & \textbf{与 Attention 的关系} & \textbf{工程含义} \\
\midrule
长上下文更容易答偏 & 相关位置增多，注意力更难集中 & 输入越长，越要重视噪声控制和检索质量 \\
显存快速上升 & 注意力矩阵和 KV Cache 变大 & 上下文预算不能只看模型参数量 \\
吞吐明显下降 & 每步需要处理更多上下文关系 & 长上下文系统必须提前做性能权衡 \\
\bottomrule
\end{tabularx}
\caption{Attention 在效果和成本上同时决定系统上限}
\label{tab:attention-cost}
\end{table}

\subsection{Attention 的工程化变体：不是改任务定义，而是改成本曲线}
旧稿里还单独展开过 MQA、GQA、FlashAttention 以及并行 Transformer Block 这一类“结构工程化”主题。它们在教材里也值得保留，因为这些设计往往不直接改变模型在做什么任务，却会显著改变\textbf{长上下文、KV Cache、显存和吞吐}的成本曲线。

\begin{table}[H]
\centering
\small
\begin{tabularx}{\textwidth}{>{\raggedright\arraybackslash}p{2.8cm}>{\raggedright\arraybackslash}p{3.1cm}>{\raggedright\arraybackslash}X}
\toprule
\textbf{工程化变体} & \textbf{主要在改什么} & \textbf{什么时候值得特别关注} \\
\midrule
MQA & 多个 query head 共享同一组 key/value，明显压缩 KV Cache & 在线推理、长输出、decode 成本高时尤其关键 \\
GQA & 让若干 query 头共享一组 key/value，在表达力和缓存成本之间折中 & 想比 MHA 更省资源，但又不想像 MQA 那样压得太狠时 \\
FlashAttention & 不改注意力数学形式，主要优化显存访问与分块计算路径 & 长上下文训练或推理吞吐上不去、HBM 带宽压力明显时 \\
并行 Transformer Block & 让 Attention 与 FFN 分支更并行地组织，缩短关键路径 & 体系结构已支持、目标明确是提训练或推理效率时 \\
\bottomrule
\end{tabularx}
\caption{Attention 工程化变体主要是在改成本，而不是改任务本身}
\label{tab:attention-engineering-variants}
\end{table}

这些变体背后的判断逻辑，其实都能回到一个问题：\textbf{瓶颈到底在算子计算、显存占用，还是内存带宽和缓存增长。} 以在线问答为例，很多时候真正拖慢系统的不是参数量本身，而是长对话和长检索上下文带来的 KV Cache 持续增长。此时 MQA 或 GQA 之类的设计价值就会变得非常直接，因为它们优先缓解的是解码阶段最常见的缓存压力。

FlashAttention 则是另一类思路。它并不试图改变“该看谁”这件事，而是更关心“看这些位置时，数据在显存和片上存储之间如何搬运”。旧稿里专门把 HBM、SRAM 和访存代价拿出来讲，其工程直觉值得保留：\textbf{很多长上下文问题不是算不出来，而是搬数据太贵。} 当团队发现注意力本身已经成为训练和推理的显存与带宽瓶颈时，这类优化就比继续堆更激进的业务层技巧更直接。

并行 Transformer Block 更适合放在“结构组织方式”理解。它提醒我们：现代大模型的优化并不总是改损失函数或改微调方法，有时只是重新安排层内计算次序，让硬件更容易吃满，让关键路径更短。但这类改动通常伴随更强的体系结构依赖，因此在工程上更像“选模型时要看懂的差异”，而不是项目中途最先动的螺丝。

\subsection{并行 Transformer Block：不是把 Attention 和 FFN 随手并起来}
旧稿里把并行 Transformer Block 单独拿出来讲过，这里值得补一层更工程化的理解。所谓“并行”，并不是说 Attention 和 FFN 可以完全互不依赖地乱序计算，而是指在某些结构设计里，两个分支都从同一个归一化后的输入出发，各自做变换，再把结果汇回残差路径。这样做的直接目标，是缩短层内关键路径，并提高硬件执行时的并行度。

这件事为什么重要？因为当模型已经很深、单层又很重时，很多成本不再只是“总 FLOPs 有多少”，而是“这些 FLOPs 能不能以更顺手的方式喂给硬件”。并行 Block 的价值，往往就体现在这里：\textbf{不一定改变任务能力边界，但可能改善单位时间内真正完成的有效计算。}

但教材里也要把边界讲清楚。并行组织不是免费午餐。它通常需要与 Pre-LN、残差设计和具体实现路径一起考虑，否则就容易出现“理论上更并行，实际上训练更难调”或“结构上改了，收益却被访存和调度抵消”的情况。所以对工程读者来说，更实用的结论不是“并行结构一定更好”，而是：\textbf{当你在比较底座模型时，要知道有些模型把效率收益写在层结构里，而不是只写在推理框架里。}

\subsection{FFN：模型看到信息之后怎样改写表示}
Attention 解决的是“看哪里”的问题，前馈网络 FFN 解决的是“看完之后怎样改写表示”的问题。它的典型形式可以写成
\[
\mathrm{FFN}(x)=W_2 \sigma(W_1x+b_1)+b_2.
\]
你可以把它理解成：Attention 把不同位置的证据先聚合回来，FFN 再对当前这个位置的内部表示做非线性变换，让它变得更适合后续层继续处理。

如果说 Attention 更像“信息读取器”，那么 FFN 更像“局部加工器”。它不直接跨位置搜索证据，但它决定了当前 token 经过一层之后，会不会把“这是一个审批条件”“这是一个数值限制”“这是一个例外条款”这类抽象模式编码得更清楚。

这也是为什么很多工程现象不能只看 Attention。模型并不只是“看见了就会用”。它还要在内部表示里把这些信息转换成可继续传播的结构，而 FFN 正是在做这件事。

\subsection{FFN 为什么经常是参数和算力大户}
很多初学者会天然把 Attention 当成 Transformer 里最“重”的部件，因为它最显眼，也最像模型在“思考”。但如果只看单层参数量，FFN 往往更大。设隐藏维度为 \(h\)，标准自注意力里的 \(q\)、\(k\)、\(v\) 和 \(o\) 四个投影矩阵参数量合起来大约是 \(4h^2\)；而传统 FFN 若把中间维度设成 \(4h\)，则参数量约为
\[
h\times 4h+4h\times h=8h^2.
\]
也就是说，在很多主流解码器里，FFN 才是单层里更大的参数和 FLOPs 承担者。Attention 更容易在长上下文里暴露为显存与带宽瓶颈，FFN 则更像一笔“常驻税收”：只要层数足够多、批次足够大，它就会持续吃掉大量算力预算。

这条判断对工程很有用。因为当团队讨论“要不要换激活函数”“中间维度要不要再放大”“为什么某个模型同样参数量却更慢”时，真正受影响的往往不只是一个局部实现细节，而是整层 FFN 的预算分配。很多模型技术报告里那些看上去不太整齐的中间维度，本质上都是在这里做精算。

\begin{note}
如果系统已经能检索到正确条款，却总在最后整理结论、合并条件或补全字段时出错，先别把责任全推给检索层。很多时候，是 FFN 没有把已经取回的证据稳定改写成可继续传播的当前位置表示。
\end{note}

\subsection{为什么现代大模型偏爱门控激活}
在现代大模型里，SwiGLU、GEGLU 这类门控激活往往优于传统 ReLU。原因不只是“论文效果更好”，而是它们在表达能力和训练稳定性之间常能取得更好的折中。以 SwiGLU 为例，可以近似写成
\[
\mathrm{SwiGLU}(x)=\mathrm{Swish}(xW)\otimes(xV),
\]
其中 \(\otimes\) 表示逐元素乘法。它比普通激活更像“带闸门的表示变换器”。

工程上可以把门控激活理解成两层作用：
\begin{itemize}
  \item 一层在做非线性变换，提升表达能力。
  \item 另一层在做选择，决定哪些信号应该被保留、哪些应该被抑制。
\end{itemize}

这类结构为什么重要？因为真实任务中，模型不只要“会变换”，还要“会克制”。企业知识助手中常见的格式遵循、字段输出、结构化调用，背后都依赖模型能稳定地保留某些模式而抑制另一些模式。门控激活在这类任务里通常更有优势。

\subsection{SwiGLU 的中间维度为什么常写成 \texorpdfstring{$\frac{8}{3}h$}{8/3h}}
很多教材第一次讲到 LLaMA、Qwen 或 Mistral 一类模型时，读者会困惑：为什么 FFN 中间维度经常不是整齐的 \(4h\)，而是 11008、14336 这类数？答案通常不是“作者偏好”，而是参数预算算出来的。

传统 FFN 有两个大矩阵，所以当中间维度取 \(4h\) 时，总参数量约为 \(8h^2\)。门控 FFN 则有三组线性变换，若把中间维度记作 \(d_{\mathrm{ffn}}\)，参数量近似是 \(3h\cdot d_{\mathrm{ffn}}\)。为了让门控 FFN 的预算与传统 FFN 保持同一量级，就会令
\[
3h\cdot d_{\mathrm{ffn}}\approx 8h^2,
\qquad
d_{\mathrm{ffn}}\approx \frac{8}{3}h.
\]
这就是为什么很多模型的门控 FFN 中间维度看起来像是“\(4h\) 的三分之二”。

以 \(h=4096\) 为例，\(\frac{8}{3}h\approx 10922.7\)，工程上再向硬件友好的倍数取整，就很容易出现 11008 这样的配置。这个数字不是装饰，它意味着：\textbf{既想保留门控激活带来的表达优势，又不想让 FFN 参数量和计算量直接膨胀 50\%}。反过来说，如果你把 SwiGLU 直接套在 \(4h\) 的中间维度上，模型当然可能更强，但代价会立刻体现在显存、吞吐和部署成本里。

这也解释了为什么有些模型在引入 MQA 或 GQA 后，会把省下来的预算重新分配给 FFN 中间维度或隐藏维度。模型设计并不是“这里省一点、那里随便花一点”，而是在不同瓶颈之间重新分配同一笔总预算。

\subsection{MoE：不是让每次都算更多，而是让总参数量和单次算量脱钩}
旧稿在长文本与模型结构部分还提过混合专家（MoE）。它值得在这里保留，因为它补了一种现代模型设计思路：\textbf{并不是所有层都必须让所有参数同时参与一次前向。} 在 MoE 结构里，路由器会根据当前 token 的表示，只激活少数几个专家网络，而不是让整块稠密 FFN 全量工作。

把它写成一个最小公式，可以记成：
\[
\mathrm{MoE}(x)=\sum_{i \in \mathrm{Top}\text{-}k} g_i(x)\,E_i(x),
\]
其中 \(E_i\) 是第 \(i\) 个专家，\(g_i(x)\) 是路由器给出的门控权重。工程直觉很简单：模型可以拥有更大的总参数池，但每个 token 前向时只支付其中一部分计算。

\begin{table}[H]
\centering
\small
\begin{tabularx}{\textwidth}{>{\raggedright\arraybackslash}p{2.8cm}>{\raggedright\arraybackslash}p{3cm}>{\raggedright\arraybackslash}X}
\toprule
\textbf{结构路线} & \textbf{主要优点} & \textbf{最该警惕的工程代价} \\
\midrule
稠密 FFN & 路径稳定、实现成熟、训练心智简单 & 总参数一大，单次前向和训练成本也一起上升 \\
MoE / 稀疏专家 & 总参数量可以更大，单次激活计算不必同步膨胀 & 路由不均衡、专家负载偏斜、并行与通信复杂度更高 \\
\bottomrule
\end{tabularx}
\caption{MoE 的核心不是“更多参数”，而是“参数总量与单次激活成本不再强绑定”}
\label{tab:moe-intuition}
\end{table}

这条机制对工程读者很重要，因为它解释了为什么一些现代模型会呈现出“总参数很大，但单次推理成本没按同样比例上涨”的现象。但教材里也必须讲清楚它的边界：MoE 不是免费午餐。只要路由不均衡、专家长期吃不到样本、跨设备专家分发太重，训练和部署复杂度就会上升。所以从原理上，你只需要先抓住一句话：\textbf{MoE 是在改“哪些参数该被激活”，而不是在改 Attention 和 FFN 解决的问题本身。}

\subsection{残差路径：为什么信息不能每层都推倒重来}
Transformer 之所以能做得很深，除了 Attention 和 FFN，本质上还依赖残差连接。残差路径允许模型在每一层都保留一条“原信息直通”的通道，而不是要求每层都从零重新构造表示。

工程上可以把残差理解成两层含义：
\begin{itemize}
  \item 它让梯度更容易在深层网络里传播，降低训练难度。
  \item 它让每一层更像是在“增量修正”表示，而不是“彻底重写”表示。
\end{itemize}

这对企业问答很有帮助。因为很多时候，模型需要保留先前层已经识别出的关键信号，只在后面逐渐补充条件、限定范围和语言组织。如果每一层都把前一层的信息覆盖得太狠，模型会更容易丢掉关键约束。

\subsection{DeepNorm 与残差缩放：模型越深，残差越不能随便加}
旧稿在 Norm 章节里专门提过 DeepNorm，这里值得保留，因为它背后讲的是一个很实际的问题：\textbf{模型层数一深，残差更新的幅度如果不受控，训练就会开始变得脆。} 残差连接让信息可以直通，但“直通”不等于“每层都可以毫无节制地往上叠加同量级的改动”。

DeepNorm 一类方案的直觉，就是通过残差缩放和初始化协同设计，让每层更新既能持续起作用，又不至于在很多层叠加后把表示尺度推爆。工程上可以把它理解成一种“深层预算管理”：每层不是不能改，而是每层都要在一个更可控的幅度里改。

这件事对教材尤其重要，因为它能帮助读者理解一个常见现象：为什么同样是 Transformer，浅一点的模型用普通残差和 LayerNorm 也能训，层数一深却开始需要更讲究的 Norm 和残差配方。答案往往不在某个单独公式，而在\textbf{深层网络里所有小更新叠加后的总效应。} 这也是为什么当项目后面进入继续预训练、对齐或大规模训练时，结构上的“稳定性设计”会重新变得非常值钱。

\subsection{Norm：模型如何保持训练稳定}
LayerNorm 的常见形式为
\[
\mathrm{LN}(x)=\gamma \cdot \frac{x-\mu}{\sqrt{\sigma^2+\epsilon}}+\beta.
\]
它的作用不是让数字“更好看”，而是控制不同层之间的表示尺度，使深层网络在前向和反向传播时更稳定。没有 Norm，深层 Transformer 很容易在训练中出现梯度不均衡、数值波动和收敛困难。

从工程视角看，Norm 解决的是“信息在很多层之间传播时，能不能不失控”。模型越深、训练越大、微调越激进，这个问题就越关键。后面讲微调和继续预训练时，我们还会再看到这一点：很多训练稳定性问题，本质上都绕不开奖励、学习率和 Norm 的共同作用。

\begin{note}
当你在训练日志里看到 loss 大幅振荡、不同层更新幅度不均、梯度异常或者同样超参数在两个模型上表现差很多时，Norm 设计通常是最值得优先排查的对象之一。
\end{note}

\subsection{RMSNorm、Pre-LN 与位置设计的工程取舍}
现代大模型里，一个很重要的趋势是大量采用 RMSNorm 和 Pre-LN 结构。你不需要先记住所有变体公式，但至少要知道三点：
\begin{itemize}
  \item \textbf{RMSNorm} 通常比标准 LayerNorm 更简洁，开销更低，又能保留足够稳定性，因此在现代解码器模型里很常见。
  \item \textbf{Pre-LN} 通常比 Post-LN 更适合深层模型，因为它在训练时往往更稳。
  \item \textbf{位置设计} 与 Norm 不是完全独立的。长上下文能力能不能稳定，不只看位置编码本身，也看整层结构如何传播这些位置信号。
\end{itemize}

这背后共同说明的一件事是：现代模型结构并不是一味追求“更复杂”，而是在效果、稳定性和硬件效率之间不断折中。后面你看到不同模型的结构差异时，最好默认它们背后都有明确的工程代价函数，而不是作者的审美选择。

\subsection{Norm 方案怎么选：先问你更怕什么}
如果把 Norm 方案选择翻译成工程语言，它通常不是“哪个公式更优美”，而是“你当前更怕什么问题”。有的团队更怕深层训练直接不稳，有的更怕推理和训练成本太高，有的则是在做超深模型时更怕残差更新失控。不同方案的优先目标并不一样。

\begin{table}[H]
\centering
\small
\begin{tabularx}{\textwidth}{>{\raggedright\arraybackslash}p{2.6cm}>{\raggedright\arraybackslash}p{3cm}>{\raggedright\arraybackslash}X}
\toprule
\textbf{方案} & \textbf{优先解决什么} & \textbf{工程上的典型代价或提醒} \\
\midrule
Post-LN + LayerNorm & 浅层或中等深度模型里的性能基线 & 层数一深，梯度更容易不稳，不适合作为现代深层解码器的默认起点 \\
Pre-LN + LayerNorm & 先把深层训练稳定性保住 & 往往更稳，但不是所有场景都追求同一类上限表现 \\
Pre-LN + RMSNorm & 在保持稳定的同时降低归一化开销 & 现代 Decoder-Only 模型常用默认项，但仍需配合学习率、初始化和批大小一起看 \\
DeepNorm 一类残差缩放方案 & 超深层、超大规模训练时抑制残差更新爆炸 & 需要连同初始化一起设计，更像体系级方案，而不是中途随手可换的小插件 \\
Sandwich-LN 一类额外归一化方案 & 对激活值爆炸有额外顾虑的特殊场景 & 实现更重，也可能引入新的训练不稳定，不适合作为大多数项目默认项 \\
\bottomrule
\end{tabularx}
\caption{Norm 方案的差异主要体现在“先解决什么问题”}
\label{tab:norm-choice}
\end{table}

对工程团队最有用的记忆法是：\textbf{Pre-LN 决定你能不能稳地把梯度送下去，RMSNorm 决定这种稳定能否以更低成本维持，DeepNorm 决定模型继续加深时残差更新会不会失控。} 当你发现同一套学习率、warmup 和批大小在不同底座模型上手感完全不同，先查这些结构习惯，通常比先怀疑“数据是不是脏了”更高效。

\subsection{交叉熵：模型训练到底在优化什么}
自回归预训练最常见的目标，是最小化交叉熵损失：
\[
\mathcal{L}_{\mathrm{CE}}=-\sum_{t=1}^{T}\log p_\theta(x_t \mid x_{<t}).
\]
它衡量的是模型对标准答案 token 的概率分配是否足够集中。直观地说，交叉熵越低，意味着模型越倾向于把正确 token 放到更高概率位置。

这里要特别注意一个常见误区：交叉熵优化的是“更像训练分布中的正确答案”，而不是“对业务最有帮助”。如果训练数据里的答案普遍不带引用、不强调拒答、风格模糊，那么模型即使把交叉熵降得很好，也未必会变成一个更适合企业知识助手的系统。

这就是为什么后面会有 SFT、对齐训练和评测体系。因为单纯把语言模型目标做到更好，并不等于把业务目标做到更好。

\subsection{一次预训练更新：损失信号怎样穿过整个网络}
如果把本章前面的部件串起来看，一次预训练更新其实没有那么神秘。模型先读入一串真实前文 token，把它们变成嵌入表示；再经过位置设计、Attention、FFN、Norm 和残差层层传播；最后输出一个词表分布，并在“正确下一个 token”上计算交叉熵。到这里，前向过程结束，反向过程才真正告诉模型哪里做错了。

这个误差信号会先作用在输出头上，要求它把正确 token 的 logit 提高、把不该高的候选压低；然后梯度继续往回流，经由最后几层的 FFN 和 Attention，把“为什么这一步没有把正确 token 放到前列”的责任逐层分摊回去；再往前，它还会回到输入嵌入和更早层的表示，使模型慢慢学会在类似上下文里把更有用的信息传到更高层。

这条反向路径的工程价值在于，它解释了为什么看似只发生在“最后一层”的预测错误，会反过来塑造整个网络的习惯。训练不是只在教输出层背词，而是在持续调整：哪些位置该被关注，哪些模式该被放大，哪些表示更有利于把正确 token 推到分布前列。所以一旦训练数据带着错误引用、模糊边界或不稳定格式，这些坏信号也会沿着同样路径被写回模型内部。

\subsection{Teacher Forcing：为什么训练看起来很顺，生成时却仍可能跑偏}
自回归训练里还有一个特别关键、但初学时容易被忽略的事实：训练阶段模型通常是在\textbf{真实前文}上预测下一个 token，而推理阶段模型却要在\textbf{自己刚刚生成出来的前文}上继续往后写。这种差异常被称为 teacher forcing 与实际生成之间的缝隙。

它带来的工程后果很直接。训练时，模型即使前一步本来容易写错，也仍然会被“正确前文”扶回正轨，所以 loss 可能看起来很健康；可一旦到了真实生成阶段，只要早期某一步偏了，后面看到的上下文就已经不是训练时那条干净轨迹，错误很可能一路放大，最终表现成答非所问、格式漂移或引用链断裂。

这条机制也解释了为什么有些模型在离线 loss 上看起来进步明显，真实长回答却未必同步变稳。因为交叉熵优化的是“给定真实前文时，下一个 token 排名是否更对”，而不是“生成整段长答案时，每一步自举前文都还能不跑偏”。后面讲采样、提示约束和对齐时，这个缝隙会继续反复出现。

\subsection{为什么分类与生成任务通常不用 MSE}
旧稿里专门讨论过“分类为什么不用均方误差”，这件事在教材里也值得留一层工程直觉。MSE 更像是在数值空间里比较“离得有多远”，而交叉熵更像是在概率分布上比较“正确答案有没有被排到前面”。对分类和生成来说，我们真正在乎的是后者。

这会带来两个直接差别。第一，当模型对错误答案已经非常自信时，MSE 往往给不出足够强的纠偏信号，而交叉熵仍会对“自信但错误”的预测保持很强惩罚。第二，语言模型面对的是一个巨大词表，它的任务不是把输出拟合成某个连续数值，而是把正确 token 尽量推到分布前列。因此，交叉熵与 Softmax 的组合更贴近任务本身。

这条判断的价值在于，它提醒我们：训练目标不仅要“能算”，还要和任务的决策结构匹配。语言模型要解决的是概率排序问题，而不是回归问题。

\subsection{Softmax 与数值稳定性：概率归一化不是实现细节}
旧稿在损失函数部分还专门讲过 Softmax 的数值稳定性，这一层在教材里也值得保留。因为模型并不是先得到“好概率”，再顺便拿去算交叉熵，而是先得到一组 logits，再通过 Softmax 变成分布。这个过程中如果数值范围失控，训练会直接变得不稳定。

最常见的问题来自指数放大。某些 logit 特别大时，直接做指数运算很容易溢出；某些 logit 特别小时，又可能让有效梯度过弱。于是工程实现里常见的减最大值技巧，本质上不是编程小聪明，而是在保证“概率归一化”这件事本身能可靠发生。

这件事为什么值得工程读者记住？因为它会反过来影响很多你在训练日志里看到的现象。例如：loss 异常、概率过尖、梯度不稳定、某些类别几乎不再被分到概率质量，背后都可能和 logits 尺度或归一化过程有关。后面讲采样时我们还会再次遇到这件事，因为推理里的温度调节，本质上也是在改 Softmax 前的 logit 形状。

\subsection{熵与信息增益：为什么有时先澄清一句更值钱}
旧稿里还有一节“信息增益”，放到教材主线里并不需要展开成完整信息论推导，但它很值得保留一层工程直觉。若一个分布越分散，说明模型越不确定；若某个额外信息能显著缩小这种不确定性，它就带来了更高的信息增益。对企业知识助手来说，这个想法并不抽象，因为它直接对应一个常见动作：\textbf{先问一个澄清问题，再回答。}

例如用户问“这个能报吗”，模型如果不知道“这个”指的是交通、住宿还是餐饮，就可能在多个制度条款之间摇摆。此时最有价值的动作，不一定是勉强给出一个模糊答案，而是先补一个能最大幅度减少不确定性的条件，比如“请说明报销类型”或“请补充岗位与差旅场景”。从信息增益角度看，这其实是在主动请求最能压缩后验不确定性的证据。

把这个视角放进教材里，有两个好处。第一，它能帮助读者理解：模型不是“答得出来就该立即答”，而是要看当前分布是否已经足够集中。第二，它也把后面第 4 章的上下文管理和澄清策略提前埋了地基。很多高质量交互，并不是因为模型一开始就知道全部事实，而是因为它能识别“现在最缺哪一条信息，补上后最能减少歧义”。

\subsection{KL：为什么对齐训练总绕不开它}
KL 散度通常写成
\[
D_{\mathrm{KL}}(p\|q)=\sum_i p(i)\log\frac{p(i)}{q(i)}.
\]
对工程团队来说，它最重要的意义不是数学定义，而是：\textbf{KL 用来约束新策略不要偏离参考分布太远。} 在对齐训练、蒸馏和策略优化里，KL 常被用来防止模型为了追逐某个局部奖励而把原本的语言能力和常识行为一起破坏掉。

如果把交叉熵理解成“朝着标准答案学”，那可以把 KL 理解成“别学偏太远”。二者在现代大模型训练中经常配合出现：前者推动模型靠近任务目标，后者防止模型失去原有基础能力。

\subsection{交叉熵与 KL 的关系：谁在决定模型“更像谁”}
把交叉熵和 KL 分开记并不难，更关键的是看懂它们之间的关系。若目标分布记为 \(p\)，模型分布记为 \(q\)，则有
\[
H(p,q)=H(p)+D_{\mathrm{KL}}(p\|q).
\]
当目标分布 \(p\) 固定时，\(H(p)\) 是常数，所以最小化交叉熵，本质上等价于让模型分布去逼近目标分布。工程上这句话非常重要，因为它能把很多看似不同的训练阶段放进同一个框架里理解。

预训练是在逼近语料里的下一个 token 分布，SFT 是在逼近示范回答分布，蒸馏是在逼近教师模型分布，而对齐训练则是在“逼近新偏好目标”的同时，再用 KL 把当前策略锚回参考模型附近。于是你会发现，很多训练方法表面差异很大，底层其实都在回答同一个问题：\textbf{我希望模型更像谁，又允许它偏离谁多远。}

一旦这样理解，很多工程决策就会清楚很多。比如“要不要加 KL 约束”“蒸馏时教师模型该选谁”“SFT 样本风格要不要统一”，本质上都不是孤立技巧，而是在改目标分布或改允许偏离的范围。

\subsection{多任务 loss 不平衡：同一个模型学多件事时谁在主导更新}
旧稿里还讨论过多任务学习时不同 loss 尺度差异过大的问题。教材里更适合把它翻译成一个直白判断：\textbf{当模型同时学多件事时，真正决定更新方向的，往往不是“你想让它学什么”，而是“哪个 loss 在优化器眼里更大声”。}

这在大模型工程里非常常见。一个系统可能同时追求语言流畅、结构化输出、工具调用成功率、检索对齐或安全拒答。如果只是把几个目标简单相加，某个数值更大、梯度更猛的目标就可能长期压制其他目标。最后得到的并不是“全面更强”的模型，而是“某一维被过度训练、其他维被挤压”的模型。

因此，多任务训练除了“有哪些目标”，还必须关心“这些目标怎样一起进入更新”。常见控制手段包括：调权重、分阶段训练、做 curriculum、先训基础任务再训高约束任务，或者直接把某些目标放回系统层而不是硬写进参数。这个判断很关键，因为它能帮你理解：为什么有时 loss 总和在下降，系统行为却在变差。

\subsection{相似度与对比学习：为什么后面检索训练看起来不像生成训练}
旧稿在损失函数后面还讲了相似度和对比学习，这部分虽然不属于生成模型本体，却是后面 RAG 检索训练的重要地基。生成任务常用交叉熵，是因为它在学“下一个 token 该是什么”；而检索器常做的是另一件事：\textbf{让正确的 query-doc 对更近，让错误的对更远。}

于是，工程里会频繁看到余弦相似度、点积和对比损失。它们回答的不是“该生成哪个词”，而是“哪个片段更像真正支持答案的证据”。对工程读者来说，最重要的不是背公式，而是记住这层差别：同样叫训练，生成模型和检索模型追逐的目标结构并不一样。

这也解释了为什么本书后面讲 RAG 优化时，会反复回到正例、负例、难负例和对比学习。它们不是突然闯进来的新世界，而是当前这一章里“损失定义决定模型行为”这条原则在检索器上的延伸。尤其要注意，\textbf{负样本太容易时，loss 也许会很好看，但检索边界未必真的被学锋利。} 真正能逼模型学会区分近似条款、相邻制度或相似问法的，往往是那些“看起来很像，但其实不能支持答案”的难负例。

\subsection{负样本为什么决定检索边界}
如果把对比学习讲得再落地一点，会得到一个很重要的工程判断：\textbf{检索训练的难度，往往不在正样本，而在负样本。} 正样本通常比较好找，比如“这个问题的正确支持文档就是这段制度条款”；真正难的是确定哪些“看起来很像，但其实不该被召回”的文本该拿来当负例。

负样本太容易时，模型学到的只是粗糙区分，例如把“差旅报销制度”和“服务器采购流程”分开。这种训练虽然 loss 漂亮，却很难帮助系统在真实环境里区分“差旅住宿上限”和“差旅交通补贴”这类高相似文本。相反，难负例越接近真实混淆边界，模型越有机会把判别面学锋利。

但难负例也有一个常见陷阱：\textbf{伪负例。} 某些文本虽然不是标注中的“唯一标准答案”，却实际上也能支持问题的一部分。如果把这类文本简单打成负例，模型会被教得过于激进，后面在 RAG 中反而容易漏召回。因此，负样本构造不是单纯“找不相干文本”，而是要在“足够难”和“别把正确证据误伤”之间平衡。这也是为什么后面讲检索优化时，我们会强调人工抽查、交叉验证和多来源证据，而不是把负例挖掘完全交给自动规则。

\subsection{为什么分类、生成和对齐都会绕回这些损失}
对工程实践而言，交叉熵和 KL 不是只会在公式区出现的名词，而是训练现象的解释器：
\begin{itemize}
  \item 在预训练里，交叉熵决定模型是否更像语料那样继续生成。
  \item 在 SFT 里，交叉熵决定模型是否更像示范答案那样组织输出。
  \item 在对齐训练里，KL 决定新策略能否在变得更“听话”的同时，不把原有语言能力一起带坏。
\end{itemize}

这直接指向一个很重要的工程结论：\textbf{loss 不是一个抽象分数，而是优化器真正追逐的目标。} 如果目标定义错了，训练越稳定，模型可能越偏。

\begin{table}[H]
\centering
\small
\begin{tabularx}{\textwidth}{>{\raggedright\arraybackslash}p{2.7cm}>{\raggedright\arraybackslash}p{3.1cm}>{\raggedright\arraybackslash}X}
\toprule
\textbf{训练场景} & \textbf{主要损失} & \textbf{工程含义} \\
\midrule
预训练 & 交叉熵 & 学会像训练语料那样继续生成 \\
SFT & 交叉熵 & 学会像任务示范那样组织回答 \\
对齐与策略优化 & 奖励目标 + KL & 在优化偏好的同时不偏离原模型太远 \\
\bottomrule
\end{tabularx}
\caption{同样的模型，不同训练阶段在追逐不同目标}
\label{tab:loss-roles}
\end{table}

\subsection{从症状反推机制}
只学公式而不会反推症状，原理就很难变成工程能力。下面给出一组最常见的现象到机制映射：
\begin{itemize}
  \item \textbf{长上下文时更容易答偏}：优先回看 Attention、位置设计和无关信息比例。
  \item \textbf{格式经常漂移}：优先看训练目标、提示模板和任务数据，而不只是怀疑解码参数。
  \item \textbf{微调过程不稳定}：优先看 Norm、学习率、批大小和梯度行为。
  \item \textbf{训练 loss 降了，但业务效果没涨}：优先看交叉熵目标和真实业务目标之间是否脱节。
  \item \textbf{显存突然吃紧}：除参数量外，也要优先怀疑 Attention 结构和 KV Cache。
\end{itemize}

这张“症状到机制”的映射表，会在后面很多章节里反复出现。因为真正的工程学习不是背下更多名词，而是遇到问题时，先知道要回到哪一层。

\section{主案例推进}
回到企业知识助手，这些机制分别解释了系统中的典型现象。

\subsection{为什么检索片段一多，回答反而更差}
很多团队第一次做 RAG 时，会觉得“把更多片段喂给模型，总不会更差”。但一旦理解了 Attention，就会知道事情没有这么简单。检索片段越多，相关位置越多、噪声也越多，模型的注意力越可能被分散。于是，明明正确证据已经在上下文里，最终回答却仍然跑偏。

这个现象说明：RAG 的问题并不只发生在检索层，也会在模型的上下文选择层暴露出来。后面我们讲 RAG 时，会反复强调“控制噪声”和“维持证据密度”，本质就是在帮 Attention 减负。

\subsection{为什么把表格摊平成纯文本后更容易漏字段}
企业知识助手里经常会遇到费用标准表、审批矩阵、岗位权限表。对人来说，表格天然带有行列结构；对语言模型来说，如果这些内容被粗暴摊平成纯文本，它看到的却只是一长串 token 顺序。此时它必须依靠位置关系、分隔符和局部模式，自己把“列名”“适用对象”“数值阈值”“例外条件”重新拼出来。

一旦分块切得不好、列间边界不清、上下文里又混入了其他制度段落，Attention 很难稳定对齐“字段名对应哪一个值”，FFN 也更难把这些关系改写成清晰表示。于是最终现象就会变成：模型好像每个词都见过，但字段一多就开始漏项、串列或把例外条件接错对象。

这正是为什么后面讲文档理解和检索优化时，本书不会把 PDF 解析、表格结构化和检索召回分开看。对系统来说，版式损失最终会变成上下文理解损失；对模型来说，输入结构一旦塌了，后面的 Attention 和 FFN 只能在更差的底座上工作。

\subsection{为什么结构化输出不只是提示词问题}
企业知识助手经常要求模型输出固定字段、审批步骤、引用来源或工具调用参数。很多人一开始会把格式问题全部当成提示词问题，但从本章视角看，格式稳定性至少同时受三层因素影响：
\begin{itemize}
  \item Attention 是否看到了真正关键的格式约束。
  \item FFN 是否能把这些约束稳定地编码进当前位置表示。
  \item 训练目标和样本是否真的强化了这种结构化输出模式。
\end{itemize}

所以，当系统总在字段顺序或引用格式上漂移时，不要只想到“提示词再写严一点”。有时问题更深，是模型没有在相应分布上被充分训练过。

\subsection{为什么有些微调一开始就不稳}
企业知识助手后面一定会接触 SFT、LoRA 和继续预训练。到了那一步，大家经常会问：为什么同样的学习率，在模型 A 上很稳，在模型 B 上却频繁 loss spike？从本章视角看，这往往与 Norm、残差路径和目标分布变化共同有关。

也就是说，微调不是“给一点数据，跑一遍训练”这么简单。底座模型的结构习惯、层间稳定性和原始分布，都会决定它对新任务的适应方式。理解这一点，后面看微调路线时就不会把所有训练不稳都误判成“数据不够多”。

真到训练现场时，一个更实用的排查顺序是：先看 warmup 初期的梯度是否异常抖动，再看 spike 是否集中在长样本或高损失样本上，最后再比较同一批数据在另一个结构更稳的底座上是否明显更容易收敛。前两步帮助你判断问题更像数值稳定与结构传导，最后一步帮助你区分“底座习惯差异”和“数据本身就有问题”。

\section{常见坑与工程取舍}
\begin{itemize}
  \item \textbf{只关注参数量，不关注上下文机制}：更大的模型不一定更适合长文档，Attention 和位置设计同样关键。
  \item \textbf{把 FFN 和激活函数当次要细节}：在同等预算下，激活函数会明显影响吞吐、稳定性和最终效果。
  \item \textbf{忽略残差和 Norm 的训练角色}：很多训练稳定性差异，并不是“玄学”，而是结构设计差异。
  \item \textbf{把交叉熵下降当成业务成功}：loss 降了，不代表系统更会引用、更会拒答或更符合企业规范。
  \item \textbf{误用 KL 约束}：KL 太强会让新策略几乎学不到东西，太弱又会让模型行为漂移失控。
\end{itemize}

再补三类工程上特别常见的取舍：
\begin{itemize}
  \item \textbf{长窗口还是高证据密度}：很多场景下，缩短噪声比一味拉长窗口更重要。
  \item \textbf{更复杂结构还是更稳的训练}：现代模型大量结构折中，目标往往是效果与稳定性的联合最优，而不是单点极致。
  \item \textbf{更低 loss 还是更强业务行为}：训练目标必须服务于系统目标，而不是只服务于日志里的分数。
\end{itemize}

\section{本章练习}
\begin{exercise}[长上下文答偏]
企业知识助手在长制度文档问答中，明明相关句子已经被放进提示，但回答仍然经常引用前文无关会议纪要。这个现象更可能首先对应哪个机制？应该优先缩短什么内容？
\end{exercise}

\begin{solution}
更可能首先对应{\heiti Attention} 的问题。模型在过长上下文中把注意力分散到了无关片段上。应优先缩短无关历史消息和低相关检索片段，而不是一开始就怀疑模型完全不懂制度内容。
\end{solution}

\begin{exercise}[损失与业务目标]
某次 SFT 后训练 loss 明显下降，但模型仍经常不给引用，只直接输出结论。这说明了什么？
\end{exercise}

\begin{solution}
这说明{\heiti 训练目标与业务目标之间可能存在脱节}。交叉熵下降只说明模型更像训练示范答案，不说明它一定学会了“先给依据再给结论”。应先检查训练样本是否真的把引用写成必选输出结构，而不是只看 loss。
\end{solution}

\section{延伸阅读}
\begin{itemize}
  \item \textbf{Vaswani et al., \emph{Attention Is All You Need}}：理解 Transformer 最初设计意图的起点。
  \item \textbf{Llama、Qwen、Mistral 等模型技术报告}：适合理解现代解码器在 Norm、激活、位置设计上的实际工程取舍。
  \item \textbf{Sebastian Raschka, \emph{Build a Large Language Model (From Scratch)}}：适合把公式、实现和训练目标对起来看。
  \item \textbf{FlashAttention、MQA、GQA 相关资料}：适合理解为什么结构变化会直接影响推理成本和长上下文部署。
\end{itemize}

\section{小结}
本章给出了工程上最值得掌握的一条 Transformer 数据流：模型先把文本变成带位置信息的表示，再通过 Attention 决定看哪里，通过 FFN 决定怎样改写表示，在残差和 Norm 的帮助下稳定传播，最后在交叉熵和 KL 等目标下被训练。理解这条数据流后，企业知识助手中的很多症状就不再是模糊现象，而能被映射回具体机制。

\section{下一章过渡}
理解底层机制后，下一步就应该把它们接到真实代码与推理链路里。下一章将从 Transformers 库出发，建立从“加载模型”到“读输出、调采样、做调试”的最小使用闭环。
